% ======================================================================
\section{Conclusions}\label{sec:conclusions}
% ======================================================================
% Sources: all sections; sec_dglap_rc.tex (notes/fix_poles_dglap.tex, ref3_beta_scheme.tex, ref3_collinear.tex),
% sec_finite.tex (notes/fix_soft.tex, fix_singlet.tex, ref3_finiteness.tex, ref3_prescriptions.tex, ref3_cubic.tex),
% sec_checks.tex (ref3_dilute.tex, ref3_cubic.tex, ref3_liskokov.tex).

\purpose{We summarise the results and then list, completely, what is not settled.}

We have computed the single inclusive gluon spectrum at mid rapidity in the collision of a dilute projectile with a dense target to next-to-leading order in the Color Glass Condensate, in pure Yang--Mills theory. The calculation uses the evolution operator $\Omega$ of the projectile, written as a normal-ordered series in the soft gluon operators with coefficients that are functionals of the non-commuting valence charge density, and reduces the spectrum to products of these coefficients and of Wilson lines. We have organised the $O(g^4)$ cross section into three groups of terms and fifteen rows. The colour charges of the projectile have been treated as non-commuting operators throughout; their ordering, the contact terms left by their commutators and the symmetrised three-point function of the projectile are part of the result. The complete result, with pointers to all its terms, is summarised in Section~\ref{sec:result_summary}. Rapidity evolution of projectile and target, collinear evolution of the final state, the running of the coupling and the finite terms all emerge from this one calculation, which in this sense completes the NLO description of the single inclusive gluon spectrum within the eikonal, dilute-projectile approximation, up to the open points listed in Section~\ref{sec:openpoints}.

The rapidity-enhanced terms have the structure required by high-energy factorisation. The logarithm of the longitudinal cut-offs splits at the momentum of the produced gluon. The coefficient of $\log(\vee/k^+)$ is the BFKL evolution of the projectile charge correlator inserted into the LO spectrum, and the coefficient of $\log(k^+/\Lambda)$ is the JIMWLK Hamiltonian acting on the target operator of the LO spectrum, for an arbitrary ordered projectile correlator. Terms with four charges cancel, as required by the coherent-state structure of the LO wave function, and the terms with three charges cancel in the symmetrised sense; with non-commuting charges, their reordering produces part of the two-charge structures of both evolutions. That no rapidity logarithm can multiply the C-odd three-point function of the projectile follows from charge conjugation together with the factorised form of the logarithms; the explicit cancellation confirms it. The operator ordering of the charges, in particular of the three-charge term of the $O(g^3)$ wave function, is essential for these cancellations.

All poles $1/\epsilon$ of the calculation are accounted for, in conventional dimensional regularisation with the $\MS$ coupling and the $\MS$ fragmentation function. The ultraviolet poles of the vertex correction of the observed gluon and of the term in which the two gluons of the two-gluon component recombine cancel each other, together with their $\frac{11}{3}$ terms; this is a result, while the finite terms that accompany it depend on the grouping of the two rows: with the colour of the gluon pair evaluated at its centre of mass, the self-energy part of the vertex correction, including its constant $\frac{67}{9}-\frac{\pi^2}{3}$, is cancelled identically. The coupling is renormalised by the ultraviolet half of the normalisation of the dressed produced gluon, and the infrared half of that normalisation completes the real final-state splitting to the full $\MS$ splitting function $\mathcal P_{gg}$, whose collinear pole is absorbed into the fragmentation function. A left-over pole multiplied by $\log(1-k^+/\vee)=O(k^+/\vee)$ is a power correction outside the eikonal approximation, which we define as the leading term of the expansion in $k^+/\vee$ and $\Lambda/k^+$ at fixed $\epsilon$. The collinear term is the final-state DGLAP fragmentation of the produced gluon. Its end points $\xi\to1$ and $\xi\to0$ are subtracted into the JIMWLK and the BFKL term, respectively, which is the ``double plus'' prescription of~\cite{Kovner:2023vsy}. The fragmentation function is the standard $\MS$ number density; with this convention the one-hadron-inclusive gluon coefficient functions of $e^+e^-\to gX$ and of $H\to gX$ are reproduced exactly. The argument of the collinear logarithm is then the transverse separation of the observed gluon, conjugate to $\v k$, and the form of~\cite{Chirilli:2012jd} corresponds to fragmentation functions that differ from the standard ones by a finite, scale-independent convolution, which we give. The running-coupling term is the $\frac{\beta_0}{2}\delta(1-\xi)$ part of the collinear term, added and subtracted; with separate renormalisation and factorisation scales it leaves $\frac{\alpha_s\beta_0}{4\pi}\log(\mu_R^2/\mu_F^2)\times$LO. With it, the coupling of the LO spectrum runs with the transverse separation conjugate to the produced momentum or, in the local form, with the produced momentum itself. Which of the numbers $\frac{11}{6}$ of the calculation are $\beta_0/(2N_c)$ is identified by a BLM-type quark loop; the coefficient of the large-distance tail of the collinear row is not.

The finite terms are given Group by Group and Row by Row, with the exact cancellations carried out and every product of three charges written as its symmetrised product plus two-charge contact terms, for an arbitrary projectile and, after the average over a globally colour-invariant projectile, in colour-singlet form with traces of adjoint Wilson lines. All colour structures are written through differences of Wilson lines, as at leading order, and all longitudinal integrals with plus prescriptions in a momentum fraction. They contain no renormalisation scale and no rapidity cut-off. The soft remainders of the plus prescriptions in the rows in which the observed gluon recoils against a soft gluon cancel in pairs of a virtual and a real term with the same coefficient, colour structure and kernel, without any condition on the target at large distance; equivalently, they can be absorbed into a JIMWLK term in which the soft gluon is restricted to rapidities below that of the observed gluon. All kernels, including those of the contact terms, are integrable at all coincidences; the contact terms do not cancel among the rows, and they do not need to. The measured-momentum Fourier integrals are unique for any regulator of the measured gluon, and with them done first all remaining integrals converge absolutely in the grouping we give, except for principal values and one conditionally convergent collinear point, which is defined by the rotation-invariant regularisation in which its subtraction was obtained. The large-distance tail of the collinear row cancels against two other rows, including the region in which only part of the gluon pair is far. The only energy denominators that vanish inside the integration region occur in the final-state dressing of the observed gluon, in one row. They are defined as principal values. The instantaneous terms of the two-gluon interference, which converge for $\vee\to\infty$ only as distributions, come from $p^+\lesssim k^+\max\{1,1/(|\v k|\ell)\}$, $\ell$ the shortest transverse scale of the integrand, not from the cut-off, and they do not depend on the shape of the separation of soft and valence modes. In the averaged form the terms cubic in the charges contribute in two ways, both given row by row. Their symmetrised part is proportional to the symmetrised three-point function $\sigma$ of the projectile, for which no model is assumed; it is odd in $\v k$ configuration by configuration, and its colour factors are C-odd target structures with at most three Wilson lines, the four-line invariants cancelling between two rows. The commutators of the ordered charges leave contact terms of the same order as the quadratic terms, which are even in $\v k$ and present for any projectile. The split of each cubic row into these two parts is unique: under charge conjugation it is the split into the C-odd and C-even parts of the projectile, as we verified on a single quark and on a dipole computed exactly in charge space.

In the dilute--dilute limit the real emission reproduces exactly the two-gluon production cross section of quasi-multi-Regge kinematics~\cite{Fadin:1998py,Bartels:2006hg}, including its normalisation and its polarisation structure, and its finite part coincides with the corresponding remainder of~\cite{Bartels:2006hg}. For a point-like projectile the virtual correction has the gluon Regge trajectories as coefficients of the two rapidity logarithms and contains no $\frac{67}{9}$; in this limit the result does not depend on the prescription for the vanishing energy denominators.

\subsection{Open points}\label{sec:openpoints}

The following points are not settled by this paper. Table~\ref{tab:open} gives an overview; the complete statements follow.
\begin{table}[t]
\centering
\renewcommand{\arraystretch}{1.25}
\footnotesize
\begin{tabular}{@{}l>{\raggedright\arraybackslash}p{0.30\textwidth}>{\raggedright\arraybackslash}p{0.56\textwidth}@{}}
\toprule
 & open point & what would settle it\\
\midrule
(i) & dilute limit of the virtual correction & real and virtual parts with matched regulators; comparison with the one-loop central-emission vertex\\
(ii) & parity in $\v k$ of the cubic terms against Li and Skokov & origin of the relative phase $i$ (prescription for $1/k^+$); kernel-level comparison in momentum space\\
(iii) & prescription for the vanishing energy denominators & Row~I.III with an $i\varepsilon$ prescription\\
(iv) & numerical evaluation & momentum-space form of all finite terms; integration with realistic Wilson lines\\
(v) & constant and $\beta_0$-type terms left in Row~I.V & evaluation of the genuine two-$\rho$ part of Row~I.V; target model\\
(vi) & collinear limit & comparison with $pp\to gX$ at NLO; $d_\perp$ continuation at the collinear point of Row~I.V\\
(vii) & status of the integrability statements & proofs for degenerate configurations; further numerical checks\\
\bottomrule
\end{tabular}
\caption{Overview of the open points; the complete statements are in the text.}
\label{tab:open}
\end{table}

\begin{enumerate}
\item[(i)] \emph{Dilute limit.} The comparison of the virtual correction with the one-loop central-emission vertex~\cite{Fadin:1998py,DelDuca:1998cx,Bartels:2006hg} has not been completed. It requires virtual and real parts with matched regulators; in particular, it is open whether the target-side term $\frac{\alpha_sN_c}{\pi}[\frac2\epsilon\log\frac{k^+}{\Lambda}+\log^2\frac{k^+}{\Lambda}]$ of the virtual correction for a point-like quark, Eq.~\eqref{eq:dilute_groupI}, cancels against the real emission from the region $|\v p|\lesssim(\Lambda/k^+)|\v k|$, which we have not computed, and whether the finite function $C(\v k,\v l)$ agrees with the one-loop vertex up to terms that can be absorbed into the unintegrated gluon distributions. Also not done are the conversion of the virtual trajectories between our plus-momentum cut-offs and the energy scales of~\cite{Bartels:2006hg}, the reduction of the virtual correction for a general projectile, the dilute limit of the coordinate-space finite terms as printed, and the comparison of the printed colour-singlet contact terms with the commutator contributions required in this limit.
\item[(ii)] \emph{Comparison with Li and Skokov.} Our symmetrised cubic terms and the first saturation correction of~\cite{Li:2021ntm} agree in their colour structures but not in their parity in $\v k$: their result has a $\v k$-even part, which our symmetrised terms do not have, and in the two-Wilson-line sector their even term has no counterpart here, while our odd two-line term has none there. We traced the difference to a relative phase $i$ in the two-charge structures that build the non-linear Weizs\"acker--Williams field and in the rescattering on the second charge after the collision; its origin, possibly the prescription for $1/k^+$ (retarded or symmetric $x^-$ profile), has not been determined, and the even part of~\cite{Li:2021ntm} has not been checked independently. A kernel-level comparison of the odd three-line terms, and the question whether our odd two-line terms add up to zero, are open.
\item[(iii)] \emph{Prescription.} The finite terms are given for the principal-value definition of the vanishing energy denominators of the final-state dressing of the observed gluon only (Section~\ref{sec:fin_conv}); the dependence of the result on this definition has not been studied here.
\item[(iv)] \emph{Numerical evaluation.} No numerical evaluation of the finite terms, and no plot, is given. Apart from the dilute limit, the finite terms are given in coordinate space only; a momentum-space form of all finite terms, and a complete numerical integration over all transverse positions with realistic Wilson lines, have not been carried out.
\item[(v)] \emph{Constant and $\beta_0$ terms.} The LO-proportional constant that is left in the finite terms of the recombination row in our grouping has not been evaluated for a general projectile and target, and the momentum-weighted collinear regions of three rows (II.I, III.III, I.V) have not been integrated. The logarithm of the emitter distance times the inverse saturation scale, with a $\beta_0$-type coefficient, that the resolving region of the recombination row contains by power counting, has not been evaluated with a target model. The $N_f$ terms are not part of this pure Yang--Mills calculation; the quark loop is used only for the identification of the $\beta_0$ terms.
\item[(vi)] \emph{Collinear limit.} The collinear convention is fixed by the general $\MS$ counterterm and by two angle-integrated coefficient functions; a comparison with the NLO coefficient of hadroproduction, $pp\to gX$, and a collinear limit of the hybrid formula have not been carried out. At the conditionally convergent collinear point of the recombination row, a continuation of the transverse indices to $d_\perp=2-\epsilon$ could give a finite local term; whether it is present has not been decided.
\item[(vii)] \emph{Status of the integrability statements.} The power counting of the contact kernels is numerical (exponent fits at generic configurations, two smooth target fields) together with analytic arguments, not a proof for every degenerate configuration; the kernel of the $f^{abc}\{\rho^b,\rho^c\}$ term of the vertex correction was estimated analytically only. The absence of a contribution from the region in which only part of the gluon pair is far was verified numerically for one charge and a target uniform at large distance; for a colour-neutral projectile and a non-uniform target it rests on an argument. The interchange of the $p^+$ integral with the joint large-distance limit was checked numerically for one of the two rows concerned. The statements on the instantaneous terms were checked with model functions of the correct short-distance behaviour, and the universality of the shift of the BFKL logarithm under a change of the soft--valence separation is argued from the form of the integrands, not checked row by row.
\end{enumerate}
